\documentclass[11pt,a4paper]{article}

\usepackage[utf8]{inputenc}
\usepackage[T1]{fontenc}
\usepackage[brazil]{babel}
\usepackage{lmodern}
\usepackage{geometry}
\geometry{margin=2.4cm}
\usepackage{graphicx}
\usepackage{float}
\usepackage{booktabs}
\usepackage{longtable}
\usepackage{array}
\usepackage{amsmath,amssymb}
\usepackage{caption}
\usepackage{subcaption}
\usepackage{xcolor}
\usepackage{hyperref}
\usepackage{siunitx}
\sisetup{output-decimal-marker={,}}

\title{Projeto 2: Quantização de Cores com SOM e GNG}
\author{Gabriely Di Folco Rocha\\
\small Introdução às Redes Neurais Artificiais}
\date{14 de outubro de 2026}

% Ajuste estes caminhos se o relatório estiver em outra pasta.
% Estrutura esperada: relatório na raiz de Project2/, imagens em figures/resized/
% e reconstruções em results/.
\newcommand{\placeholder}[1]{%
  \fbox{\parbox[c][4.0cm][c]{0.88\linewidth}{\centering #1}}}
\newcommand{\TEpct}[1]{#1\%}

\begin{document}
\maketitle

\section{Procedimento Experimental}

\subsection{Objetivo e imagens}

Este projeto compara duas redes neurais competitivas não supervisionadas,
Self-Organizing Map (SOM) e Growing Neural Gas (GNG), na quantização de cores
de imagens. O objetivo é representar a variedade de cores de cada imagem com
um conjunto limitado de protótipos RGB e avaliar tanto a fidelidade da
representação quanto uma medida de organização topológica.

Cada pixel é tratado como uma amostra tridimensional,
\[
\mathbf{x}=(R,G,B)\in[0,1]^3.
\]
Os valores dos canais são normalizados dividindo-se os valores originais por
255. A posição espacial do pixel não é usada no treinamento: a rede recebe
somente os três valores de cor. Após o treinamento, cada pixel é substituído
pelo vetor de pesos do neurônio mais próximo, denominado vencedor ou
\textit{Best Matching Unit} (BMU). A imagem reconstruída mantém as dimensões
originais, mas utiliza somente as cores dos protótipos aprendidos.

Foram selecionadas três imagens para a comparação principal: \texttt{berlim},
\texttt{canudos} e \texttt{debret11}. As imagens foram redimensionadas para que
o maior lado tivesse no máximo 500 pixels, preservando-se a proporção. A
Tabela~\ref{tab:images} apresenta o número de pixels utilizado na avaliação.

\begin{table}[H]
\centering
\caption{Imagens selecionadas para os experimentos.}
\label{tab:images}
\begin{tabular}{lrr}
\toprule
Imagem & Pixels avaliados & Pixels na amostra de treino\\
\midrule
Berlim & 178\,500 & 5\,000\\
Canudos & 185\,500 & 5\,000\\
Debret 1 & 150\,000 & 5\,000\\
\bottomrule
\end{tabular}
\end{table}

\begin{figure}[H]
\centering
\begin{subfigure}[t]{0.31\linewidth}
\centering
\IfFileExists{figures/resized/berlim.png}
{\includegraphics[width=\linewidth]{figures/resized/berlim.png}}
{\placeholder{Inserir imagem original\\Berlim}}
\caption{Berlim}
\end{subfigure}\hfill
\begin{subfigure}[t]{0.31\linewidth}
\centering
\IfFileExists{figures/resized/canudos.png}
{\includegraphics[width=\linewidth]{figures/resized/canudos.png}}
{\placeholder{Inserir imagem original\\Canudos}}
\caption{Canudos}
\end{subfigure}\hfill
\begin{subfigure}[t]{0.31\linewidth}
\centering
\IfFileExists{figures/resized/debret11.png}
{\includegraphics[width=\linewidth]{figures/resized/debret1.png}}
{\placeholder{Inserir imagem original\\Debret 1}}
\caption{Debret 1}
\end{subfigure}
\caption{Imagens utilizadas na comparação.}
\label{fig:originals}
\end{figure}

\subsection{Redes e treinamento}

O SOM organiza seus neurônios em uma grade bidimensional fixa. Em cada
apresentação, identifica-se o BMU e atualizam-se seus pesos e os pesos dos
neurônios vizinhos, com influência determinada pela distância na grade. Foram
testadas grades \(4\times4\), \(8\times8\) e \(16\times16\), correspondentes a
16, 64 e 256 neurônios.

O GNG começa com dois neurônios e modifica a topologia durante o treinamento.
O algoritmo atualiza o vencedor e seus vizinhos, mantém conexões entre
neurônios, remove conexões antigas e insere novos neurônios em regiões de erro
acumulado elevado. Foram testados limites de 16, 64 e 256 neurônios. O número
final de neurônios observado em cada execução atingiu o limite configurado.

Para cada imagem, foi selecionada uma amostra de 5\,000 pixels sem reposição.
A mesma amostra foi reutilizada entre configurações e sementes para permitir
comparações controladas. Cada modelo foi treinado com as sementes 42, 43 e 44,
totalizando três execuções independentes por configuração. O orçamento de
treinamento do GNG foi de 60\,000 passos, amostrando pixels com reposição. Para
o SOM, o orçamento foi aproximado por 12 épocas completas sobre a amostra de
5\,000 pixels, também totalizando 60\,000 apresentações de pixels.

Os parâmetros usados no GNG foram: taxa de aprendizagem do vencedor igual a
0,05, taxa dos vizinhos igual a 0,0006, idade máxima das arestas igual a 50,
inserção de neurônios a cada 100 passos, \(\alpha=0,5\) e decaimento do erro
acumulado por fator \(1-\beta\), com \(\beta=0,0005\). Esses valores descrevem
a implementação utilizada neste experimento.

No SOM, a taxa de aprendizagem e o raio de vizinhança foram reduzidos
linearmente ao longo das épocas. A implementação utilizada leva esses valores
próximos de zero na última época. Essa programação difere da configuração da
referência consultada, na qual a taxa de aprendizagem termina em um valor
positivo e o raio de vizinhança termina em 0,5; portanto, os resultados não
devem ser interpretados como uma reprodução exata daquela configuração.

\begin{table}[H]
\centering
\caption{Configurações da comparação principal.}
\label{tab:params}
\begin{tabular}{p{0.34\linewidth}p{0.56\linewidth}}
\toprule
Parâmetro & Configuração\\
\midrule
Imagens & Berlim, Canudos e Debret 1\\
Amostra por imagem & 5\,000 pixels, sem reposição\\
Sementes & 42, 43 e 44\\
Orçamento de treinamento & 60\,000 apresentações/passos por execução\\
SOM & Grades \(4\times4\), \(8\times8\) e \(16\times16\)\\
GNG & Limites de 16, 64 e 256 neurônios\\
GNG: taxa do vencedor & 0,05\\
GNG: taxa dos vizinhos & 0,0006\\
GNG: idade máxima de aresta & 50 passos\\
GNG: inserção & A cada 100 passos\\
GNG: parâmetros de erro & \(\alpha=0,5\), \(\beta=0,0005\)\\
\bottomrule
\end{tabular}
\end{table}

\subsection{Métricas}

\paragraph{Erro de quantização (QE).}
O QE é a distância euclidiana RGB média entre cada pixel avaliado e seu
protótipo vencedor. Para \(M\) pixels avaliados,
\[
\mathrm{QE}=\frac{1}{M}\sum_{j=1}^{M}
\left\|\mathbf{x}_j-\mathbf{w}_{b(j)}\right\|_2,
\]
em que \(b(j)\) é o índice do BMU do pixel \(\mathbf{x}_j\) e
\(\mathbf{w}_{b(j)}\) é o vetor de pesos correspondente. Valores menores
indicam que as cores reconstruídas estão, em média, mais próximas das cores
originais.

\paragraph{Erro topológico (TE).}
O TE é a fração de pixels para os quais o primeiro e o segundo neurônios mais
próximos não são vizinhos na topologia da rede. No SOM, são vizinhos os
neurônios adjacentes horizontal ou verticalmente na grade; no GNG, são
vizinhos os neurônios ligados por uma aresta no grafo final. Se \(a_1(j)\) e
\(a_2(j)\) são os dois neurônios mais próximos do pixel \(j\), então
\[
\mathrm{TE}=\frac{1}{M}\sum_{j=1}^{M}
\mathbb{1}\!\left[a_1(j)\not\sim a_2(j)\right].
\]
O TE é apresentado em porcentagem nas tabelas. Como a topologia do SOM é fixa
e a do GNG é aprendida, os valores devem ser interpretados conforme a
estrutura de cada modelo. Um TE menor não implica, isoladamente, melhor
qualidade global de quantização.

Para cada combinação de imagem, modelo e configuração, foram calculadas a
média e o desvio-padrão amostral das três sementes. O desvio-padrão descreve a
variação entre execuções, não a incerteza associada a uma amostra aleatória de
todas as imagens possíveis.

\section{Resultados quantitativos}

As Tabelas~\ref{tab:berlim}--\ref{tab:debret1} apresentam média \(\pm\) desvio
padrão amostral. O QE é apresentado na escala RGB normalizada; o TE é
apresentado em porcentagem.

\begin{table}[H]
\centering
\caption{Resultados para Berlim.}
\label{tab:berlim}
\small
\begin{tabular}{llrr}
\toprule
Modelo & Configuração & QE médio \(\pm\) DP & TE (\%) médio \(\pm\) DP\\
\midrule
SOM & \(4\times4\) & 0,05347 \(\pm\) 0,00074 & 49,01 \(\pm\) 5,15\\
SOM & \(8\times8\) & 0,03222 \(\pm\) 0,00029 & 53,60 \(\pm\) 0,39\\
SOM & \(16\times16\) & 0,02260 \(\pm\) 0,00012 & 57,75 \(\pm\) 0,85\\
GNG & 16 neurônios & 0,05335 \(\pm\) 0,00042 & 0,47 \(\pm\) 0,62\\
GNG & 64 neurônios & 0,03093 \(\pm\) 0,00011 & 0,76 \(\pm\) 0,19\\
GNG & 256 neurônios & 0,02098 \(\pm\) 0,00008 & 3,36 \(\pm\) 0,27\\
\bottomrule
\end{tabular}
\end{table}

\begin{table}[H]
\centering
\caption{Resultados para Canudos.}
\label{tab:canudos}
\small
\begin{tabular}{llrr}
\toprule
Modelo & Configuração & QE médio \(\pm\) DP & TE (\%) médio \(\pm\) DP\\
\midrule
SOM & \(4\times4\) & 0,03013 \(\pm\) 0,00012 & 57,51 \(\pm\) 2,21\\
SOM & \(8\times8\) & 0,01456 \(\pm\) 0,00014 & 70,61 \(\pm\) 0,38\\
SOM & \(16\times16\) & 0,00943 \(\pm\) 0,00015 & 68,59 \(\pm\) 2,30\\
GNG & 16 neurônios & 0,02982 \(\pm\) 0,00008 & 0,00 \(\pm\) 0,00\\
GNG & 64 neurônios & 0,01328 \(\pm\) 0,00002 & 0,22 \(\pm\) 0,14\\
GNG & 256 neurônios & 0,00808 \(\pm\) 0,00002 & 2,04 \(\pm\) 0,16\\
\bottomrule
\end{tabular}
\end{table}

\begin{table}[H]
\centering
\caption{Resultados para Debret 1.}
\label{tab:debret1}
\small
\begin{tabular}{llrr}
\toprule
Modelo & Configuração & QE médio \(\pm\) DP & TE (\%) médio \(\pm\) DP\\
\midrule
SOM & \(4\times4\) & 0,04465 \(\pm\) 0,00006 & 57,34 \(\pm\) 7,30\\
SOM & \(8\times8\) & 0,02748 \(\pm\) 0,00012 & 60,78 \(\pm\) 3,25\\
SOM & \(16\times16\) & 0,01917 \(\pm\) 0,00010 & 56,07 \(\pm\) 1,17\\
GNG & 16 neurônios & 0,04441 \(\pm\) 0,00010 & 0,09 \(\pm\) 0,09\\
GNG & 64 neurônios & 0,02656 \(\pm\) 0,00003 & 0,70 \(\pm\) 0,43\\
GNG & 256 neurônios & 0,01775 \(\pm\) 0,00005 & 3,53 \(\pm\) 0,54\\
\bottomrule
\end{tabular}
\end{table}

\subsection{Interpretação das métricas}

Nas três imagens, o QE diminuiu quando aumentou o número de neurônios, tanto
no SOM quanto no GNG. Esse padrão é compatível com a maior capacidade de
representação oferecida por uma paleta com mais protótipos. Em configurações
com o mesmo número nominal de neurônios, o GNG apresentou QE médio ligeiramente
menor que o SOM nas três imagens. As diferenças são empíricas e estão
condicionadas às amostras, sementes e parâmetros utilizados.

Os desvios-padrão do QE foram pequenos em relação às médias, sugerindo que a
métrica foi relativamente estável entre as três sementes deste experimento.
Essa estabilidade não elimina a possibilidade de variação em outras imagens,
amostras ou configurações de treinamento.

O TE do SOM permaneceu elevado nas três imagens e não decresceu de forma
monótona com o aumento da grade. Em Canudos, por exemplo, o TE médio aumentou
de 57,51\% na grade \(4\times4\) para 70,61\% na grade \(8\times8\), antes de
diminuir para 68,59\% na grade \(16\times16\). O resultado mostra que melhorar
a quantização de cores não garante, por si só, reduzir o erro topológico.

O TE do GNG foi baixo em todas as configurações avaliadas, embora tenha
aumentado na configuração de 256 neurônios. Essa medida verifica a adjacência
entre os dois protótipos mais próximos segundo o grafo final do GNG; portanto,
não é diretamente equivalente à adjacência na grade fixa do SOM. O TE deve ser
usado como diagnóstico da organização relativa dos vencedores, e não como
medida isolada da fidelidade visual.

\section{Reconstruções das imagens}

As reconstruções são apresentadas em mosaicos para facilitar a comparação entre
as configurações. Cada linha corresponde a uma imagem e cada coluna a uma
configuração. As figuras são geradas pelo script
\texttt{src/generate\_mosaics.py} a partir das reconstruções salvas em
\texttt{results/}. Os mosaicos de paletas SOM são construídos a partir dos
pesos finais exportados pelo treinamento; por isso, é necessário acrescentar
ao laço de treinamento a gravação dos pesos em arquivo, conforme indicado no
código complementar.

\subsection{Berlim}

\begin{figure}[H]
    \centering

    \begin{minipage}{0.32\textwidth}
        \centering
        \includegraphics[width=\linewidth]{results/berlim_som_4x4_seed42.png}
        \caption*{(a) SOM $4\times4$}
    \end{minipage}
    \hfill
    \begin{minipage}{0.32\textwidth}
        \centering
        \includegraphics[width=\linewidth]{results/berlim_som_8x8_seed42.png}
        \caption*{(b) SOM $8\times8$}
    \end{minipage}
    \hfill
    \begin{minipage}{0.32\textwidth}
        \centering
        \includegraphics[width=\linewidth]{results/berlim_som_16x16_seed42.png}
        \caption*{(c) SOM $16\times16$}
    \end{minipage}

    \caption{Reconstruções de Berlim pelo SOM, com semente 42 e diferentes dimensões de grade.}
    \label{fig:berlim-som}
\end{figure}

\subsection{Berlim}

\begin{figure}[H]
    \centering

    \begin{minipage}{0.32\textwidth}
        \centering
        \includegraphics[width=\linewidth]{results/berlim_gng_16.png}
        \caption*{(a) GNG 16}
    \end{minipage}
    \hfill
    \begin{minipage}{0.32\textwidth}
        \centering
        \includegraphics[width=\linewidth]{results/berlim_gng_64.png}
        \caption*{(b) GNG 64}
    \end{minipage}
    \hfill
    \begin{minipage}{0.32\textwidth}
        \centering
        \includegraphics[width=\linewidth]{results/berlim_gng_256_seed42.png}
        \caption*{(c) GNG 256}
    \end{minipage}

    \caption{Reconstruções de Berlim pelo GNG, com semente 42 e diferentes dimensões de grade.}
    \label{fig:berlim-som}
\end{figure}

\subsection{Berlim}

\begin{figure}[H]
    \centering

    \begin{minipage}{0.32\textwidth}
        \centering
        \includegraphics[width=\linewidth]{results/palettes/berlim_som_4x4_seed42_weights_palette.png}
        \caption*{(a) paleta SOM $4\times4$}
    \end{minipage}
    \hfill
    \begin{minipage}{0.32\textwidth}
        \centering
        \includegraphics[width=\linewidth]{results/palettes/berlim_som_8x8_seed42_weights_palette.png}
        \caption*{(b) SOM $8\times8$}
    \end{minipage}
    \hfill
    \begin{minipage}{0.32\textwidth}
        \centering
        \includegraphics[width=\linewidth]{results/palettes/berlim_som_16x16_seed42_weights_palette.png}
        \caption*{(c) paleta SOM $16\times16$}
    \end{minipage}

    \caption{Paletas geradas nessa abordagem SOM.}
    \label{fig:berlim-som}
\end{figure}

\subsection{Canudos}

\begin{figure}[H]
    \centering

    \begin{minipage}{0.32\textwidth}
        \centering
        \includegraphics[width=\linewidth]{results/canudos_som_4x4_seed42.png}
        \caption*{(a) SOM $4\times4$}
    \end{minipage}
    \hfill
    \begin{minipage}{0.32\textwidth}
        \centering
        \includegraphics[width=\linewidth]{results/canudos_som_8x8_seed42.png}
        \caption*{(b) SOM $8\times8$}
    \end{minipage}
    \hfill
    \begin{minipage}{0.32\textwidth}
        \centering
        \includegraphics[width=\linewidth]{results/canudos_som_16x16_seed42.png}
        \caption*{(c) SOM $16\times16$}
    \end{minipage}

    \caption{Reconstruções de Berlim pelo SOM, com semente 42 e diferentes dimensões de grade.}
    \label{fig:canudos-som}
\end{figure}

\subsection{Canudos}

\begin{figure}[H]
    \centering

    \begin{minipage}{0.32\textwidth}
        \centering
        \includegraphics[width=\linewidth]{results/canudos_gng_16.png}
        \caption*{(a) GNG 16}
    \end{minipage}
    \hfill
    \begin{minipage}{0.32\textwidth}
        \centering
        \includegraphics[width=\linewidth]{results/canudos_gng_64.png}
        \caption*{(b) GNG 64}
    \end{minipage}
    \hfill
    \begin{minipage}{0.32\textwidth}
        \centering
        \includegraphics[width=\linewidth]{results/canudos_gng_256_seed42.png}
        \caption*{(c) GNG 256}
    \end{minipage}

    \caption{Reconstruções de Canudos pelo GNG, com semente 42 e diferentes dimensões de grade.}
    \label{fig:berlim-som}
\end{figure}

\subsection{Canudos}

\begin{figure}[H]
    \centering

    \begin{minipage}{0.32\textwidth}
        \centering
        \includegraphics[width=\linewidth]{results/palettes/canudos_som_4x4_seed42_weights_palette.png}
        \caption*{(a) paleta SOM $4\times4$}
    \end{minipage}
    \hfill
    \begin{minipage}{0.32\textwidth}
        \centering
        \includegraphics[width=\linewidth]{results/palettes/canudos_som_8x8_seed42_weights_palette.png}
        \caption*{(b) SOM $8\times8$}
    \end{minipage}
    \hfill
    \begin{minipage}{0.32\textwidth}
        \centering
        \includegraphics[width=\linewidth]{results/palettes/canudos_som_16x16_seed42_weights_palette.png}
        \caption*{(c) paleta SOM $16\times16$}
    \end{minipage}

    \caption{Paletas geradas nessa abordagem SOM.}
    \label{fig:canudos-som}
\end{figure}
\subsection{Debret}

\begin{figure}[H]
    \centering

    \begin{minipage}{0.32\textwidth}
        \centering
        \includegraphics[width=\linewidth]{results/debret1_som_4x4_seed42.png}
        \caption*{(a) SOM $4\times4$}
    \end{minipage}
    \hfill
    \begin{minipage}{0.32\textwidth}
        \centering
        \includegraphics[width=\linewidth]{results/debret1_som_8x8_seed42.png}
        \caption*{(b) SOM $8\times8$}
    \end{minipage}
    \hfill
    \begin{minipage}{0.32\textwidth}
        \centering
        \includegraphics[width=\linewidth]{results/debret1_som_16x16_seed42.png}
        \caption*{(c) SOM $16\times16$}
    \end{minipage}

    \caption{Reconstruções de Debret pelo SOM, com semente 42 e diferentes dimensões de grade.}
    \label{fig:canudos-som}
\end{figure}

\subsection{Debret}

\begin{figure}[H]
    \centering

    \begin{minipage}{0.32\textwidth}
        \centering
        \includegraphics[width=\linewidth]{results/debret1_gng_16.png}
        \caption*{(a) GNG 16}
    \end{minipage}
    \hfill
    \begin{minipage}{0.32\textwidth}
        \centering
        \includegraphics[width=\linewidth]{results/debret1_gng_64.png}
        \caption*{(b) GNG 64}
    \end{minipage}
    \hfill
    \begin{minipage}{0.32\textwidth}
        \centering
        \includegraphics[width=\linewidth]{results/debret1_gng_256_seed42.png}
        \caption*{(c) GNG 256}
    \end{minipage}

    \caption{Reconstruções de Canudos pelo GNG, com semente 42 e diferentes dimensões de grade.}
    \label{fig:debret1-som}
\end{figure}

\subsection{Debret}

\begin{figure}[H]
    \centering

    \begin{minipage}{0.32\textwidth}
        \centering
        \includegraphics[width=\linewidth]{results/palettes/debret1_som_4x4_seed42_weights_palette.png}
        \caption*{(a) paleta SOM $4\times4$}
    \end{minipage}
    \hfill
    \begin{minipage}{0.32\textwidth}
        \centering
        \includegraphics[width=\linewidth]{results/palettes/debret1_som_8x8_seed42_weights_palette.png}
        \caption*{(b) SOM $8\times8$}
    \end{minipage}
    \hfill
    \begin{minipage}{0.32\textwidth}
        \centering
        \includegraphics[width=\linewidth]{results/palettes/debret1_som_16x16_seed42_weights_palette.png}
        \caption*{(c) paleta SOM $16\times16$}
    \end{minipage}

    \caption{Paletas geradas nessa abordagem SOM.}
    \label{fig:debret1-som}
\end{figure}

A leitura conjunta das reconstruções e das paletas permite observar dois
aspectos distintos: as reconstruções mostram como a redução da paleta afeta a
imagem inteira, enquanto as paletas exibem os protótipos na organização
espacial imposta pela grade SOM. A posição dos quadrados na paleta representa
a posição do neurônio na grade, não a posição de pixels na imagem original.

\section{Reconstruções resumidas}
\subsection{Greve dos metalúrgicos}
\begin{figure}[H]
    \centering

    \begin{minipage}{0.32\textwidth}
        \centering
        \includegraphics[width=\linewidth]{figures/resized/lula.png}
        \caption*{(a) imagem original}
    \end{minipage}
    \hfill
    \begin{minipage}{0.32\textwidth}
        \centering
        \includegraphics[width=\linewidth]{results/lula_som_4x4.png}
        \caption*{(b) SOM $4\times4$}
    \end{minipage}
    \hfill
    \begin{minipage}{0.32\textwidth}
        \centering
        \includegraphics[width=\linewidth]{results/lula_gng_16.png}
        \caption*{(c) GNG}
    \end{minipage}

    \caption{visão resumida}
    \label{fig:lula-som}
\end{figure}

\subsection{Debret}
\begin{figure}[H]
    \centering

    \begin{minipage}{0.32\textwidth}
        \centering
        \includegraphics[width=\linewidth]{figures/resized/debret2.png}
        \caption*{(a) imagem original}
    \end{minipage}
    \hfill
    \begin{minipage}{0.32\textwidth}
        \centering
        \includegraphics[width=\linewidth]{results/debret2_som_4x4.png}
        \caption*{(b) SOM $4\times4$}
    \end{minipage}
    \hfill
    \begin{minipage}{0.32\textwidth}
        \centering
        \includegraphics[width=\linewidth]{results/debret2_gng_16.png}
        \caption*{(c) GNG 16}
    \end{minipage}

    \caption{visão resumida}
    \label{fig:debret2-som}
\end{figure}

\subsection{Eliseu Visconti}
\begin{figure}[H]
    \centering

    \begin{minipage}{0.32\textwidth}
        \centering
        \includegraphics[width=\linewidth]{figures/resized/eliseu_visconti.png}
        \caption*{(a) imagem original}
    \end{minipage}
    \hfill
    \begin{minipage}{0.32\textwidth}
        \centering
        \includegraphics[width=\linewidth]{results/eliseu_visconti_som_4x4.png}
        \caption*{(b) SOM $4\times4$}
    \end{minipage}
    \hfill
    \begin{minipage}{0.32\textwidth}
        \centering
        \includegraphics[width=\linewidth]{results/eliseu_visconti_gng_16.png}
        \caption*{(c) GNG 16}
    \end{minipage}

    \caption{visão resumida}
    \label{fig:eliseu_visconti-som}
\end{figure}

\subsection{Georgina de Abulquerque}
\begin{figure}[H]
    \centering

    \begin{minipage}{0.32\textwidth}
        \centering
        \includegraphics[width=\linewidth]{figures/resized/georgina.png}
        \caption*{(a) imagem original}
    \end{minipage}
    \hfill
    \begin{minipage}{0.32\textwidth}
        \centering
        \includegraphics[width=\linewidth]{results/georgina_som_4x4.png}
        \caption*{(b) SOM $4\times4$}
    \end{minipage}
    \hfill
    \begin{minipage}{0.32\textwidth}
        \centering
        \includegraphics[width=\linewidth]{results/georgina_gng_16.png}
        \caption*{(c) GNG 16}
    \end{minipage}

    \caption{visão resumida}
    \label{fig:georgina-som}
\end{figure}

\section{Limitações e trabalho futuro}

Os resultados constituem uma comparação inicial, com três imagens, uma amostra
de treinamento de 5\,000 pixels por imagem e três sementes por configuração.
Embora as sementes permitam avaliar a variação entre execuções, a quantidade
de imagens é pequena para generalizar as conclusões a todas as classes de
imagem.

Além disso, o orçamento de treinamento do SOM foi implementado como 12 épocas
completas, com redução linear da taxa de aprendizagem e do raio de vizinhança
até valores próximos de zero na última época. Essa configuração difere da
agenda de parâmetros descrita na referência usada como modelo de organização
do relatório. Os parâmetros do GNG também não são idênticos aos da referência,
em particular a taxa dos vizinhos, a idade máxima das arestas e o decaimento do
erro acumulado. Assim, a comparação apresentada é válida para as implementações
e configurações efetivamente executadas, mas não constitui reprodução exata
dos resultados de referência.

Como continuação, recomenda-se testar uma agenda de aprendizagem do SOM que
mantenha os valores finais positivos, avaliar mais imagens e incluir
visualizações diagnósticas: U-matrix, ocupação dos neurônios, distribuição dos
canais RGB e grafo final do GNG. Essas visualizações ajudariam a explicar a
organização dos protótipos e as diferenças observadas nas métricas.

\section{Conclusões}

Neste experimento, SOM e GNG foram aplicados à quantização de cores em três
imagens, usando redes com 16, 64 e 256 neurônios. O aumento da quantidade de
neurônios reduziu consistentemente o erro de quantização em ambas as redes. O
GNG apresentou QE médio ligeiramente inferior ao SOM nas comparações com o
mesmo número de neurônios, enquanto o desvio-padrão entre as três sementes foi
pequeno para essa métrica.

O erro topológico apresentou comportamentos distintos: no SOM, os valores
foram elevados e não diminuíram monotonamente com o tamanho da grade; no GNG,
foram baixos segundo a adjacência do grafo aprendido. Essa diferença precisa
ser interpretada considerando que cada modelo define sua própria noção de
vizinhança. Portanto, os resultados favorecem o GNG quanto ao QE nas
configurações avaliadas, mas não justificam uma conclusão universal de
superioridade entre os algoritmos.

\section*{Referências}

\begin{enumerate}
\item KOHONEN, T. \textit{Self-Organizing Maps}. 3. ed. Berlin: Springer, 2001.
\item FRITZKE, B. A growing neural gas network learns topologies. In:
\textit{Advances in Neural Information Processing Systems}, v. 7, 1995.
\end{enumerate}

\end{document}